\metatitle{Generalized Demazure atoms}
\metadescription{An introduction to generalized Demazure atoms, their definitions, properties, and connections to fillings and RSK.}  


\section[generalizedDemazureAtoms]{Generalized Demazure atoms}


\begin{polydata}{generalizedDemazureAtoms}
  Name   & Generalized Demazure atoms \\
  Space    & All \\
  Basis    & True \\
  Rating   & 1 \\
  Bib      & HaglundMasonRemmel2012 \\
  Year     & 2012 \\
  Keywords & fillings, rsk \\
  Category & Schur \\
  Contains & atom | HaglundMasonRemmel2012 | scope=identity basement \\
  Contains & schurS | HaglundMasonRemmel2012 | scope=reverse-identity basement \\
  PositiveIn & atom | HaglundMasonRemmel2012 \\
\end{polydata}

The family of polynomials referred to as \emph{generalized Demazure atoms}
were introduced in \cite{HaglundMasonRemmel2012} and further studied in \cite{RemmelTiefenbruck2013}.

The generalized Demazure atoms are related to the fillings in the combinatorial model for
\hyperref[macdonaldE]{non-symmetric Macdonald polynomials}, \cite{HaglundHaimanLoehr2008} and
\hyperref[macdonaldEperm]{permuted-basement Macdonald polynomials}, \cite{Alexandersson2015gbMacdonald},
but the definitions differ.
The families studied in \cite{Alexandersson2015gbMacdonald} and \cite{AlexanderssonSawhney2019}
are referred to as \emph{permuted-basement atoms} and are not the same as 
generalized Demazure atoms.
An example below shows the difference.

The generalized Demazure atoms have nice properties with respect to insertion algorithms and \hyperref[rsk]{RSK},
while the permuted-basement atoms are compatible with Demazure--Lusztig operators and representation theory.
For a fixed basement $\sigma$, Haglund--Mason--Remmel's generalized
Demazure atoms $\widehat{E}^{\sigma}_\gamma$ refine Schur functions:
\[
\schurS_\beta(\xvec)=\sum_{\lambda(\gamma)=\beta}
  \widehat{E}^{\sigma}_\gamma(\xvec).
\]
Here, $\lambda(\gamma)$ is the decreasing rearrangement of $\gamma$.
For the reverse-identity basement this sum has a single nonzero term, so
Schur functions occur as a subfamily.  For the identity basement one recovers
Mason's Demazure atoms.  The $t$-deformed basement families are recorded here
as \hyperref[permutedBasementTAtom]{permuted-basement $t$-atoms}.

We use the description of the fillings in \cite[Sect.~2]{RemmelTiefenbruck2013}.
Let $\gamma$ be a weak composition with $n$ parts and let $\sigma\in\symS_n$.
Draw $\gamma$ as a skyline diagram, with $\gamma_i$ cells in column $i$,
and add a basement row $0$ filled with $\sigma_1,\dotsc,\sigma_n$
from left to right.  Write $F(i,j)$ for the entry in column $i$ and row $j$.
A \defin{permuted basement filling}\label{permutedBasementFillingHMR} (PBF)
of shape $\gamma$ and basement $\sigma$ is a filling of the cells with
positive integers such that
\begin{itemize}
\item entries weakly decrease up each column, starting from the basement,
\item every triple of type $A$ or type $B$ is an inversion triple, and
\item (\defin{B-increasing condition}) if $i<l$ and $\gamma_i<\gamma_l$, then
  $F(i,j-1)<F(l,j)$ for all $j\geq 1$.
\end{itemize}
The generalized Demazure atom $\widehat{E}^{\sigma}_\gamma(\xvec)$ is the sum
of $\xvec^{F}$ over all PBFs $F$, where basement entries are not counted.
It vanishes unless $\gamma_i\geq\gamma_j$ whenever $i<j$ and
$\sigma_i>\sigma_j$ \cite[Sect.~2]{RemmelTiefenbruck2013}.
Compared with the fillings for permuted-basement atoms, PBFs satisfy the
extra B-increasing condition \cite[Rem.~24]{Alexandersson2015gbMacdonald}.

\begin{example}
Let $n=2$ and $\beta=(2,1)$, so $\gamma$ is $(2,1)$ or $(1,2)$.
Below, the bottom row of each diagram is the basement.
For the identity basement $\sigma=(1,2)$, each shape has exactly one PBF:
\begin{figure}
\ytableaushort{1{\none},12,12}
\ytableaushort{{\none}2,12,12}
\end{figure}
Hence $\widehat{E}^{(1,2)}_{(2,1)}(\xvec)=x_1^2x_2$ and
$\widehat{E}^{(1,2)}_{(1,2)}(\xvec)=x_1x_2^2$.  These are the Demazure atoms
of \cite{Mason2009} with the same indices, and they sum to
$\schurS_{(2,1)}(x_1,x_2)$.

For the reverse-identity basement $\sigma=(2,1)$, the shape $(2,1)$ has two
PBFs:
\begin{figure}
\ytableaushort{2{\none},21,21}
\ytableaushort{1{\none},21,21}
\end{figure}
so $\widehat{E}^{(2,1)}_{(2,1)}(\xvec)=x_1x_2^2+x_1^2x_2=\schurS_{(2,1)}(x_1,x_2)$.
The shape $(1,2)$ has no PBF.  The B-increasing condition would require
$2=F(1,0)<F(2,1)\leq F(2,0)=1$.  Hence $\widehat{E}^{(2,1)}_{(1,2)}=0$.

The \hyperref[permutedBasementAtom]{permuted-basement atom}
$\atom^{\sigma}_\alpha(\xvec)=\macdonaldE^{\sigma}_\alpha(\xvec;0,0)$
uses the same basement and shape, but has no B-increasing condition.
In the English convention of that page, with rows instead of columns, the
basement $(2,1)$ and shape $(1,2)$ have the single filling
\begin{figure}
\ytableaushort{22,111}
\end{figure}
where the first column is the basement.  Hence
$\atom^{(2,1)}_{(1,2)}(\xvec)=x_1^2x_2\neq 0$.  Similarly
$\atom^{(2,1)}_{(2,1)}(\xvec)=x_1^2x_2+x_1x_2^2$.
With this basement the permuted-basement atoms are key polynomials with
reversed index, $\atom^{(2,1)}_\alpha=\key_{\rev(\alpha)}$
\cite[Sect.~3.1]{Alexandersson2015gbMacdonald}.  Their sum
$2x_1^2x_2+x_1x_2^2$ is not a Schur polynomial.
For the identity basement, the two families agree in this example.
\end{example}

A computer check of the same comparison for $n=3$ and $\beta=(2,1,0)$ gives
the following.
For each of the six basements $\sigma$, the two families agree on the
shapes $\gamma$ with $\gamma_i\geq\gamma_j$ whenever $i<j$ and
$\sigma_i>\sigma_j$.  On every other rearrangement $\gamma$ of $(2,1,0)$, the
generalized Demazure atom vanishes and the permuted-basement atom does not.

\name{Janine LoBue Tiefenbruck} and \name{Jeffrey B. Remmel} prove a
Murnaghan--Nakayama rule for generalized Demazure atoms
\cite{RemmelTiefenbruck2013}.  Their rule expands the product of a
power-sum symmetric function with a generalized Demazure atom as a signed
sum of generalized Demazure atoms.
